Now saying that \((xa)b\) and \((xc)d\) are
defined means that, over \(S''\), \(x\) is in
a certain relatively schematically dense open
subset, which is checked on the closed fiber
of \(X_{S''}\). Thus it works.
\end{proof}

Now let \(G\) be the quotient of \(X^2\) by
\(R\) as a sheaf for the fppf topology.
We shall define a composition law on \(G\)
as follows: let \((g,g')\in G(S')\) be
represented by a section
\(((a,b),(c,d))\) of \(X^2\times X^2(S'')\),
with \(S''\to S'\) fppf-covering.
Suppose moreover that \(X\) admits a section
\(x\) over \(S''\) such that \(a(b(cx))\)
and \(x^{-1}d\) are defined, which is
permitted by Rule 1 and (1.7), and call
\(gg'\) the class in \(G(S')\) represented
by the section \((a(b(cx)),x^{-1}d)\)
of \(X^2(S'')\).

Let us verify that \(gg'\) does not depend
on the choice of the section \(x\) and
the representative \(((a,b),(c,d))\):
indeed, let \((a',b')\sim(a,b)\),
\((c',d')\sim(c,d)\), and \(x\) be such
that \(a'(b'(c'x'))\) and \(x'^{-1}d'\)
are defined.
We may suppose that all are sections over \(S''\).
One must prove that
\[
(a(b(cx)),x^{-1}d)\sim(a'(b'(c'x')),x'^{-1}d'),
\]
that is, for a suitable section
\(z\in X(S''')\), with \(S'''\to S''\)
suitably fppf-covering, one has
\[
(z(a(b(cx))))(x^{-1}d)
=(z(a'(b'(c'x'))))(x'^{-1}d').
\]
Whenever all the products are defined, one has
\[
\begin{aligned}
(z(a(b(cx))))(x^{-1}d)
 &=((za)(b(cx)))(x^{-1}d)\\
 &=(((za)b)(cx))(x^{-1}d)\\
 &=((((za)b)c)x)(x^{-1}d)\\
 &=(((za)b)c)(x(x^{-1}d))=(((za)b)c)d,
\end{aligned}
\]
and the same identities hold with the primes.
Now, by Rule 1 and (1.7), such a \(z\) exists.
One must therefore prove that
\[
(((za)b)c)d=(((za')b')c')d'.
\]
But \((za)b=(za')b'\), since
\((a,b)\sim(a',b')\) (3.3), and one
therefore has the desired equality,
since \((c,d)\sim(c',d')\).
% typo?: the source introduces x, then uses x'; kept as printed.

Consider the natural morphism \(i:X\to G\)
defined as follows: for \(a\in X(S')\),
choose \(S''\to S'\) fppf-covering
and a section \(b\in X(S'')\)
such that \(ab^{-1}\) is defined, and set
\[
i(a)=\text{class in }G\text{ of }(ab^{-1},b).
\]
One easily verifies that this class,
which is a priori in \(G(S'')\),
does not depend on the choice of \(b\),
and thus gives a well-defined element
of \(G(S')\).
The reader will take pleasure in verifying the
\begin{proposition}{3.6}\label{XVIII.3.6}
The morphism \(i\) commutes with the
composition laws of \(X\) and \(G\),
that is, if \(a,b\in X(S')\)
are sections such that \(ab\) is
defined, one has \(i(a)i(b)=i(ab)\).
\end{proposition}
The aim of this number is the following theorem:
\begin{theorem}{3.7}\label{XVIII.3.7}
\nde{If \(X\) is smooth, separated
over \(S\), faithfully flat of finite
presentation over \(S\), then \(G/S\)
is represented by a smooth \(S\)-group
scheme of finite type over \(S\).
This is Theorem 6.6.1 of the book
\emph{Néron models} by Bosch--Lütkebohmert--Raynaud,
Springer (1990).}
Let \((X,W)\) be a group germ
over \(S\), with \(X/S\) faithfully
flat, of finite presentation, and
with separated fibers without
embedded components.
Then, with the notation above, one has
\begin{enumeratei}
\item \(G\) is a sheaf of groups.
\item \(i:X\to G\) is represented by an open immersion.
\item \(G\) is representable locally
on \(S\) for the fppf topology.
\item If \(G/S\) is representable,
then it is a flat group of finite
presentation, and \(i:X\to G\)
is schematically dense relative to \(S\).
\end{enumeratei}
\end{theorem}
Note that \(G\) is evidently
characterized by properties (i),
\ldots, (iv); one may therefore
forget the explicit construction of \(G\).
The proof proceeds in several steps:
\begin{lemma}{3.8}\label{XVIII.3.8}
(i) Let \(c\) be a section of
\(X(S)\); then the morphism
from the prescheme \(c\times X\)
(which is \(S\)-isomorphic to
\(X\)) to \(G\), given by
\(c\times X\hookrightarrow X\times X\to G\),
is a monomorphism.

(ii) Let \(\{c_i\}\), \(i\in I\),
be sections of \(X(S)\),
let \(Z=\coprod_i c_i\times X\),
and call \(R'\rightrightarrows Z\)
the equivalence relation induced
on \(Z\) by the evident morphism
\(Z\to X^2\).
Then \(R'\) is a \og gluing\fg{}
of the \(c_i\times X\simeq X\).
\end{lemma}
\begin{proof}
(i) For two sections
\((c,a),(c,a')\) of \((c\times X)(S')\)
to have the same image in \(G\),
one must have \((c,a)\sim(c,a')\),
that is, \((xc)a\sim(xc)a'\)
for a suitable \(x\in X(S'')\),
with \(S''\to S'\) fppf-covering,
whence \(a=a'\) by Rule 3.
% typo?: sim rather than equality retained as printed.

(ii) Let \(c_i,c_j\) be two
sections, and consider the
birational map \(\psi_{ji}\)
from \(X\) to itself over \(S\),
given by the formula
\(x\mapsto c_j^{-1}(c_i x)\).
This is the same map as that
given by \(x\mapsto(yc_j)^{-1}((yc_i)x)\),
if \(y\in X(S)\), as one
easily sees.
Moreover, one verifies that
\(\varphi_{ji}\) is defined
for a \(b\in X(S')\) if
and only if there exists
\(b'\) such that
\((c_i,b)\sim(c_j,b')\),
and then \(b'=\varphi_{ji}(b)\).
Let \(U_{ji}\) be the
domain of definition over
\(S\) of \(\varphi_{ji}\).
It remains to prove that
this domain of definition
is universal, that is,
if \(b\in X(S'')\),
with \(S''\to S\) arbitrary,
and \(\varphi_{ji}\) is
defined at \(b\), then
\(b\in U_{ji}(S'')\).
It amounts to the same
to prove that if \(b,b'\in X(S'')\)
are such that
\((c_i,b)\simeq(c_j,b')\),
then \(b\in U_{ji}(S'')\).
We leave the verification
of this fact, analogous
to that of (3.5), to the reader.
% typo?: psi_ji is subsequently called varphi_ji in the source.
\end{proof}
\begin{lemma}{3.9}\label{XVIII.3.9}
Suppose that \(\{c_i\}\),
\(i\in I\), are sections
of \(X(S)\) such that
\(\coprod_i c_i\times X\to G\)
is surjective as a morphism
of sheaves.
Then \(G\) is representable
and flat, of finite presentation
over \(S\), and the structural
morphism \(X^2\to G\)
is flat and of finite presentation.
\end{lemma}
\begin{proof}
The fact that \(G\) is
representable is an immediate
consequence of (3.8), and
it follows that
\(\coprod_i c_i\times X\to G\)
is an open covering.
To prove that \(X^2\to G\)
is flat, it suffices to
do so locally, hence to
prove that the rational
map \(X^2\to c_i\times X\)
induces a flat morphism
on its domain of definition.
Now this rational map
is given by
\((a,b)\mapsto(c_i,((xc_i)^{-1}(xa))b)\),
with \(x\) an arbitrary
section, and, if it is
defined at \((a,b)\in X^2(S')\),
one can find \(S''\to S'\)
fppf-covering and a section
\(x\in X(S'')\) such that
\(((xc_i)^{-1}(xa))b\)
is defined.
One easily sees that
it is therefore a flat
morphism.
It follows similarly
that it is locally of
finite presentation, hence
fppf-covering.
Now, by construction, the
relation \(R\) is effective.
Thus, by (3.5), the morphism
\(X^2\to G\) becomes of
finite presentation after
the base change \(G\leftarrow X^2\),
which is fppf-covering;
hence \(X^2\to G\)
is of finite presentation.
Let us show that \(G\to S\)
is flat and of finite
presentation.
It is flat and locally
of finite presentation,
since \(G\) is covered
by the \(c_i\times X\simeq X\).
Now \(X^2/S\) is quasi-compact,
and \(X^2\to G\) is surjective.
This proves that \(G/S\)
is quasi-compact.
To prove that \(G\to S\)
is quasi-separated, note
that one has the following
cartesian diagram:
\[
\xymatrix{
R\ar[r]^\alpha\ar[d]_\gamma & X^2\times_S X^2\ar[d]^\beta\\
G\ar[r]_\Delta & G\times_S G }.
\]
One has \(\gamma\) surjective
and \(\alpha,\beta\) quasi-compact,
hence \(\beta\alpha\) is
quasi-compact, hence \(\Delta\)
is quasi-compact.
\nde{Another way to conclude
here is by faithfully flat
descent (EGA IV$_2$, 2.7.1),
since \(\beta\) is covering
for fppf.}
% typo: "G -> S et quasi-séparé" -> "G -> S est quasi-séparé".
\end{proof}
\begin{lemma}{3.10}\label{XVIII.3.10}
Let \(\{c_i\}\), \(i\in I\),
be sections of \(X(S)\).
For \(\coprod_i c_i\times X\to G\)
to be surjective as a
morphism of fppf-sheaves,
it suffices that the
following condition hold:

For every \(S'\to S\)
and \((a,b)\in X^2(S')\),
there exists an open covering
\(\{S'_\nu\}\), \(\nu\in N\),
of \(S'\), and a function
\(N\to I\), \(\nu\mapsto i(\nu)\),
such that \((c_{i(\nu)}^{-1}a)b\)
is defined over \(S'_\nu\).
\end{lemma}
\begin{proof}
Let \(S''\to S\) be arbitrary,
and \(g\in G(S'')\).
Choose \(S'\to S''\)
fppf-covering, and a
section \((a,b)\in X^2(S')\)
representing \(g\).
Take the open covering
\(\{S'_\nu\}\) of \(S'\)
which exists by the
hypothesis of the lemma.
Then, over each \(S'_\nu\),
one has
\((a,b)\sim(c_{i(\nu)},(c_{i(\nu)}^{-1}a)b)\);
hence \(g\) is represented
by a section of
\([c_{i(\nu)}\times X](S')\)
over \(S'_\nu\), which
proves surjectivity, since
the family of morphisms
\(\{S'_\nu\to S''\}\)
is fppf-covering.
\end{proof}
\begin{lemma}{3.11}\label{XVIII.3.11}
Let \(Y/S\) be a prescheme
of finite presentation,
and let \(\{a_i\}\), \(i\in I\),
be sections of \(Y(S)\).
Let \(s_0,s_1\) be points
of \(S\) such that \(s_0\)
is a specialization of
\(s_1\), and \(Y_j\) the
fiber of \(Y/S\) at \(s_j\).
Let \(C_j\) be the
closure in \(Y_j\)
of the set of points
\(\{a_i(s)\cap Y_j\}\).
Then one has \(\dim C_1\geq\dim C_0\).
\end{lemma}
\begin{proof}
It suffices to make the
verification after a
base change \(S'\to S\)
with chosen points
\(s'_0,s'_1\) such that
\(s'_j\mapsto s_j\)
and \(s'_0\) is a
specialization of \(s'_1\).
One is therefore (EGA II 7.1.4)
reduced to the case
where \(S\) is the
spectrum of a valuation
ring \(A\), \(s_0\)
the closed point of \(S\),
and \(s_1\) the generic
point of \(S\).
Now let \(V\) be the
closure of \(C_1\) in
\(Y\).
It is clear that
\(C_0\subseteq V\),
and thus the lemma
is a consequence of
the \og well-known\fg{}
fact that an irreducible
closed subprescheme \(V\)
of a prescheme \(Y/S\)
of finite presentation
satisfies
\(\dim(V\times s_1)\geq\dim(V\times s_0)\)
if \(S\) is the
spectrum of a valuation
ring and \(V\times s_1\ne\varnothing\)
(EGA IV 13.1.6).
\end{proof}
\begin{lemma}{3.12}\label{XVIII.3.12}
Suppose that \(S\) is
the spectrum of a
local ring, with closed
point \(s_0\), and
let \(\{c_i\}\), \(i\in I\),
be sections such that
the closure \(C_0\)
of the set
\(\{c_i(s)\cap X_0\}\)
in the closed fiber
\(X_0\) has dimension
equal to \(\dim X_0=n\).
Then the condition of
Lemma 3.9 is satisfied.
\end{lemma}
\begin{proof}
First note that the
fibers of \(X/S\)
all have the same
dimension \(u\), which
follows from EGA IV
12.1.1 (i), and the
fact that \(X\) has
a rational composition
law.
Lemma (3.11) therefore
implies that, for every
morphism \(S_1\to S\)
with \(S_1\) the
spectrum of a field,
the dimension of the
closure of the set
\(\{c_i\times_S S_1\}\)
in \(X_{S_1}\)
is equal to \(n\).
Let us verify the
condition of (3.10):
let \((a,b)\in X^2(S')\).
For \((c_i^{-1}a)b\)
to be defined, it is
necessary and sufficient
that \(c_i\) be
contained in a certain
open subset \(U\subseteq X_{S'}\)
schematically dense
relative to \(S'\)
(Rule 1).
One must prove that
this is true for a
suitable \(i\), locally
on \(S'\).
It therefore suffices
to treat the case
where \(S'\) is the
spectrum of a local
ring; then the fact
that \(c_i\in U(S')\)
is checked on the
closed fiber.
One is therefore
reduced to the case
\(S'=\Spec k\),
with \(k\) a field.
Now, with the above
notation, take \(S'=S_1\).
One has \(\dim C_1=\dim X_{S_1}\),
and \(U\) is relatively
schematically dense
in \(X_{S_1}\).
Thus \(U\cap C_1\ne\varnothing\),
whence \(U\cap\{c_p\times_S S_1\}\ne\varnothing\),
and one has won.
% typo?: dimension u, then n, and c_p are retained as printed.
\end{proof}

The proof of the
theorem is now easy.
First note the
following consequence
of the finiteness
assertion of Lemma (3.9):
if \(\{A_i\}\) is an
inductive system of
rings over \(S\),
if \(\underrightarrow A=\varinjlim A_i\),
and if the hypotheses
of (3.9) hold for
\(S=\Spec\underrightarrow A\),
then one can descend
the object representing
the quotient \(G\)
of \(R\rightrightarrows X^2\)
to one of the
\(S_i=\Spec A_i\),
with the finiteness
and flatness properties
stated in (3.9).
This is the customary
passage to the limit
(EGA IV 8 and 11).
It follows that, for
the proof of (iii)
and (iv) of (3.7),
one may restrict
oneself to the case
\(S=\Spec A\), with
\(A\) a strictly
local ring.
Then let \(x_0\)
be a closed point
of the closed fiber
\(X_0\) of \(X/S\).
There exists
{\renewcommand{\thefootnote}{(*)}\addtocounter{footnote}{-1}%
\footnote{Cf. the note at the bottom
of page 15, Exp. VII.}}
an extension \(A'\)
of \(A\), local,
free and finite,
and a section of
\(X'=X\times_S S'\)
passing through the
unique point \(x'_0\)
of \(X'_0\) over \(x_0\).
Note that \(X'_0\to X_0\)
is radicial, since
\(S\) is strictly
local, hence the
residue field of
\(A\) is separably
closed.
It follows that
there exists an
inductive system
\(\{A_i\}\) of
local rings, flat
and finite over \(S\),
such that, setting
\[
\underrightarrow A=\varinjlim A_i,\qquad
\underleftarrow X=X\times_S\Spec\underrightarrow A,
\]
\(\underleftarrow X\)
has a set of
sections inducing a
dense set on the
closed fiber
\(\underleftarrow X_0\).
For every closed
point \(x_0\) of
\(X_0\), take an
extension \(A(x_0)\)
such that the
corresponding base
extension of \(X\)
admits a section
\og passing through \(x_0\)\fg,
and take as inductive
system the system
of finite tensor
products of the
\(A(x_0)\).
Note that
\(\underrightarrow A\)
is local, being a
limit of local rings.
Thus, by (3.12)
and (3.10), one
has the quotient
\(G\) over
\(\Spec\underrightarrow A\),
hence over one of
the \(\Spec A_i\)
by the remarks above,
hence locally for
the fppf topology.
% typo: d'anneaux printed as d^prime anneaux is rendered as of rings.

In fact, it follows
from the constructions
that locally for the
fppf topology one
can find a finite
set of sections
\(\{c_i\}\), \(i=1,\ldots,r\),
such that \(G\) is
covered by the
\(c_i\times X\).
(ii) and (iv)
follow easily from
this fact, and we
leave the verification
of (i) to the reader.
\begin{corollary}{3.13}\label{XVIII.3.13}
The sheaf of
groups \(G\)
determined by a
group germ \((X,W)\)
over \(S\) is
representable in each
of the following situations:
\begin{enumeratei}
\item \(S\) is artinian.
\item For every local
scheme \(S'\to S\)
of a closed point
\(s\) of \(S\),
\(X_{S'}\) has a
set of sections
inducing on the
closed fiber a set
whose closure has
dimension \(\dim X_s\).
\item \(S\) is strictly
local, and \(X/S\)
is smooth.
\item There exists
\(S'\to S\)
fppf-covering such
that \(G_{S'}\)
is representable
and affine over \(S\).
\end{enumeratei}
\end{corollary}
\begin{proof}
Indeed, (ii) is
a consequence of
(3.10) and (3.12);
(iii) follows
directly from (ii)
and \og Hensel's lemma\fg;
(iv) from descent
of affine schemes;
and (i) from
descent of group
schemes, which is
possible here since
one knows that every
finite subset of a
group over a field
is contained in an
affine open subset
(Exp. VI).
\end{proof}

\begin{thebibliography}{2}
\bibitem[1]{Weil}
A. Weil, \emph{Variétés abéliennes et courbes algébriques}.
Hermann, Paris, 1948.
\bibitem[2]{Yanagihara}
H. Yanagihara, \emph{Reduction of group varieties and transformation spaces}.
Journ. Sci. Hiroshima Univ., Ser. A-I, vol. 27, no. 1, June 1963.
\end{thebibliography}
